\begin{abcliste}
\abc The magnetic force is always perpendicular to the velocity (and to the field). It changes only the direction of the velocity, not its magnitude (it does no work), so the electrons move on a $\result{\text{circle}}$.

\abc The accelerating voltage gives the electrons the kinetic energy
\begin{align}
 \frac{1}{2}m v^2 &= e U
\end{align}
and the magnetic force is the centripetal force:
\begin{align}
 e v B &= \frac{m v^2}{r} \quad\Longrightarrow\quad v = \frac{e B r}{m}
\end{align}
Inserting $v$ into the energy equation gives $\frac{1}{2}m\left(\frac{eBr}{m}\right)^2 = eU$, therefore
\begin{align}
 \frac{e}{m} &= \qmF \\
   &= \frac{2\times\U}{\left(\B\right)^2\times\left(\r\right)^2} \\
   &= \qm \approx \result{\qmP{0}{2}}
\end{align}
in good agreement with the accepted value $\SI{1.76e11}{\coulomb\per\kg}$.

\abc A higher voltage makes the electrons faster. Since $r = \frac{m v}{e B} = \frac{1}{B}\sqrt{\frac{2 m U}{e}}$, the circle gets $\result{\text{larger}}$.
\end{abcliste}
